\documentclass[11pt]{article}
\usepackage[letterpaper,margin=1in]{geometry}
\usepackage{amsmath,amssymb,amsthm}
\usepackage{booktabs,tabularx,array}
\usepackage[colorlinks=true,linkcolor=black,citecolor=black,urlcolor=blue]{hyperref}
\hypersetup{pdftitle={Clique identification in a Dirac--Haj\'os construction need not preserve criticality},
  pdfauthor={Alex Chan}}

\newtheorem{theorem}{Theorem}
\newtheorem{lemma}{Lemma}
\newcommand{\out}{N^{+}}
\newcommand{\chiG}{\chi}
\newcommand{\edges}{E}

\title{Clique identification in a Dirac--Haj\'os construction need not preserve criticality}
\author{Alex Chan\\[2pt]
  \small Harvard University and National Bureau of Economic Research\\[-1pt]
  \small \texttt{achan@hbs.edu}}
\date{\today}

\begin{document}
\maketitle
\begin{abstract}
Jensen and Toft ask whether a generalization of the Haj\'os construction, in which two $k$-critical graphs are identified along an $h$-clique before the usual edge deletion and addition, always produces a $k$-critical graph. We give a negative answer for every $k\geq5$, already with $h=3$. A $12$-vertex $4$-critical graph supplies the required coloring rigidity. For $k=5$ the resulting graph has $23$ vertices and $68$ edges, and an edge $e$ for which $\chi(D)=\chi(D-e)=5$. The finite base-graph assertions are accompanied by complete coloring certificates and an exhaustive three-coloring table.
\end{abstract}
\noindent\textbf{Keywords.} Critical graph; chromatic number; Haj\'os construction; clique identification.\qquad
\textbf{MSC 2020.} 05C15.

\section{The question and result}
A graph is \emph{$k$-critical} if it is $k$-chromatic and every proper subgraph is $(k-1)$-colorable. In Problem~11.8 of \cite{JensenToft}, two disjoint $k$-critical graphs $G_1,G_2$ have distinguished $h$-cliques $H_1,H_2$, where $1\leq h\leq k-2$. A vertex $y_i\notin H_i$ is adjacent to $x_i\in H_i$. The cliques are identified so that $x_1$ and $x_2$ coincide; then $x_iy_i$ is deleted in each copy and $y_1y_2$ is added. The question is whether the resulting graph must be $k$-critical. The case $h=1$ is the usual Haj\'os construction; the question originated in the survey of Sachs and Stiebitz~\cite{SachsStiebitz}.

\begin{theorem}\label{thm:main}
For every integer $k\geq5$ there are two $k$-critical graphs and an admissible choice $h=3$ for which the construction of Problem~11.8 produces a graph $D_k$ containing an edge $01$ such that
\[
\chiG(D_k)=\chiG(D_k-01)=k.
\]
In particular, $D_k$ is not $k$-critical. It has $2k+13$ vertices and $k^2+15k-32$ edges; $D_5$ has $23$ vertices and $68$ edges.
\end{theorem}

\section{A finite four-critical base}
Let $B$ have vertex set $\{0,1,\ldots,11\}$. Table~\ref{tab:graph} specifies its edge set without reference to a drawing: $N^+(v)$ lists the neighbors of $v$ with larger labels. Thus $|E(B)|=24$, $H=\{0,1,2\}$ is a triangle, and vertex $3$ is adjacent to $0$ but to neither $1$ nor $2$. We abbreviate an edge $\{u,v\}$ by $uv$.

\begin{table}[ht]
\centering
\begin{tabular}{@{}c l @ {\hspace{2.3em}} c l@{}}
\toprule
$v$ & $N^+(v)$ & $v$ & $N^+(v)$ \\
\midrule
0 & $1,2,3,4,11$ & 6 & $9,10$ \\
1 & $2,6,7,10$ & 7 & $8$ \\
2 & $9,11$ & 8 & $10$ \\
3 & $5,8,11$ & 9 & $\varnothing$ \\
4 & $6,7,8,9$ & 10 & $\varnothing$ \\
5 & $6,8$ & 11 & $\varnothing$ \\
\bottomrule
\end{tabular}
\caption{The full edge specification of $B$.}\label{tab:graph}
\end{table}

\begin{lemma}\label{lem:base}
The graph $B$ is $4$-critical. Moreover, every proper $3$-coloring $c$ of $B-03-01$ satisfies
\begin{equation}\label{eq:rigidity}
 c(3)=\begin{cases}c(0),&c(0)\ne c(1),\\c(2),&c(0)=c(1).\end{cases}
\end{equation}
\end{lemma}
\begin{proof}
The edges $02$ and $12$ remain after the two deletions. If $c(0)\ne c(1)$, relabel the colors so $(c(0),c(1),c(2))=(0,1,2)$; if $c(0)=c(1)$, use $(0,0,1)$. In the first normalization the complete list of colorings is
\[
012011202001,\qquad 012022001121,
\]
and in the second it is the single coloring
\[
001122110022.
\]
The digit in position $v$, counting from zero, is $c(v)$. Hence the list gives \eqref{eq:rigidity}. Its completeness follows by extending a coloring of $0,1,2$ successively to vertices $3,4,\ldots,11$, rejecting a proposed color when it equals the color of an earlier adjacent vertex. For reproducibility, the numbers of retained partial colorings after each step are
\[
\begin{array}{c|rrrrrrrrr}
\text{last vertex}&3&4&5&6&7&8&9&10&11\\\hline
(0,1,2)&3&6&12&12&20&7&4&3&2\\
(0,0,1)&3&6&12&8&8&2&2&2&1
\end{array}
\]
with adjacency determined by Table~\ref{tab:graph}. This recurrence checks at most $3^9$ complete assignments for either normalization.

The first two strings are exactly the normalized $3$-colorings of $B-03$, since $01$ is present in that graph. Both give $c(0)=c(3)$, so restoring $03$ shows that $B$ is not $3$-colorable. Recoloring vertex $3$ with a new color in either string yields a $4$-coloring of $B$. Appendix~\ref{app:cert} provides, for each of the $24$ edges $e$, a proper $3$-coloring of $B-e$. Every vertex has an incident edge, so these certificates prove that every proper subgraph is $3$-colorable. Thus $B$ is $4$-critical.
\end{proof}

\section{Proof of the theorem}
Put $r=k-4\geq1$ and let $G=K_r\vee B$, the join of $B$ with an $r$-clique. Then $G$ is $k$-critical. Indeed, $\chiG(G)=r+4$. Deleting an edge within $B$ permits $r$ colors on the clique and three on $B$; deleting an edge within the clique permits $r-1$ colors there and four on $B$. Finally, if an edge between a clique vertex $q$ and $v\in B$ is deleted, color $B-v$ with three colors, give $v$ the color of $q$, and use $r$ colors on the clique. Each edge deletion therefore admits a $(k-1)$-coloring.

Take two disjoint copies $G_1,G_2$ of $G$. Identify their copies of $H=\{0,1,2\}$ labelwise, delete $03$ in both copies, and add the edge $3^{(1)}3^{(2)}$. Call the result $D_k$. The identified triangle is an allowed $h$-clique, since $h=3\leq k-2$, and the selected endpoints are $x_i=0$ and $y_i=3^{(i)}$. The edge $01$ remains in $D_k$.

Suppose $D_k-01$ has a $(k-1)$-coloring. The $r$-clique of each $G_i$ uses $r$ colors disjoint from the colors on its corresponding copy of $B$. Consequently each copy of $B-03-01$ is colored with at most three colors. Apply Lemma~\ref{lem:base} to both copies. If the shared vertices $0$ and $1$ have different colors, both $3^{(i)}$ have the color of $0$; if they have the same color, both $3^{(i)}$ have the color of $2$. Either case contradicts the edge $3^{(1)}3^{(2)}$. Thus $\chiG(D_k-01)\geq k$.

For the reverse bound, color $B-3$ using colors $1,2,3$; such a coloring follows, for example, by restricting the first coloring of $B-03$ displayed above. Use this same coloring on the two copies of $B-3$ so their common triangle agrees. Give $3^{(1)}$ color $4$, and color the $r$-clique of $G_1$ with $5,\ldots,k$. Give $3^{(2)}$ color $5$, and color the $r$-clique of $G_2$ with $4,6,\ldots,k$. This is a proper $k$-coloring of $D_k$. Since $D_k-01$ is a subgraph of $D_k$, both chromatic numbers equal $k$.

Each input has $12+r$ vertices and $24+12r+\binom r2$ edges. Identifying the two triangles merges three vertices and counts their three edges once; the two deletions and one addition change the edge count by $-1$. Hence $|V(D_k)|=2(12+r)-3=2k+13$ and
\[
|E(D_k)|=2\left(24+12r+\binom r2\right)-4=k^2+15k-32.
\]
This completes the proof of Theorem~\ref{thm:main}.\qed

\section*{Reproducibility}
The accompanying Python checker tests the finite certificates and sample outputs for $k=5,6,7$. The Lean 4/mathlib file \texttt{D5Counterexample.lean} checks the $23$-vertex output: edge $01$ is present, the graph has a proper five-coloring, and deleting $01$ still permits no four-coloring. The symbolic family is defined by \texttt{GeneralConstructionNative.lean}, and \texttt{GeneralNoKMinusOne.lean} proves that its output admits no coloring with $k-1$ colors for every $k\geq5$, using the base-coloring rigidity lemma. The general $k$-coloring and the $k$-criticality of both input graphs have not yet been formalized in Lean. The equality linking the symbolic graph at $k=5$ with the explicit edge list uses \texttt{native\_decide}, which relies on Lean's compiler for that finite computation. Code and verification files: \url{https://github.com/alexgocardinal/dirac-hajos-jensen-toft-11-8}.

\appendix
\section{Edge-deletion certificates for $B$}\label{app:cert}
For each edge $e$ listed below, the corresponding twelve-digit string is a proper $3$-coloring of $B-e$, with its $v$th digit assigned to vertex $v$. Every edge other than $e$ has differently colored ends, while the endpoints of $e$ have the same digit.
\begin{center}
\small
\begin{tabular}{@{}c l @ {\hspace{2.4em}} c l@{}}
\toprule
$e$ & coloring of $B-e$ & $e$ & coloring of $B-e$\\
\midrule
$01$ & \texttt{001122110022} & $02$ & \texttt{010112020222}\\
$03$ & \texttt{012011202001} & $04$ & \texttt{012200221101}\\
$0\,11$ & \texttt{012110202000} & $12$ & \texttt{011112020222}\\
$16$ & \texttt{012220101001} & $17$ & \texttt{012221010121}\\
$1\,10$ & \texttt{012211220011} & $29$ & \texttt{012211020221}\\
$2\,11$ & \texttt{012110202002} & $35$ & \texttt{012222001121}\\
$38$ & \texttt{012210202001} & $3\,11$ & \texttt{012110202001}\\
$46$ & \texttt{012220201001} & $47$ & \texttt{012221020121}\\
$48$ & \texttt{012210201001} & $49$ & \texttt{012211020121}\\
$56$ & \texttt{012220001121} & $58$ & \texttt{012221001121}\\
$69$ & \texttt{012211020021} & $6\,10$ & \texttt{012211220021}\\
$78$ & \texttt{012221000121} & $8\,10$ & \texttt{012211220001}\\
\bottomrule
\end{tabular}
\end{center}

\begin{thebibliography}{9}
\bibitem{JensenToft}
T. R. Jensen and B. Toft, \emph{Graph Coloring Problems}, Wiley, New York, 1995, \S 11.8, pp. 185--186. \href{https://doi.org/10.1002/9781118032497.ch11}{doi:10.1002/9781118032497.ch11}.
\bibitem{SachsStiebitz}
H. Sachs and M. Stiebitz, Colour-critical graphs with vertices of low valency, in \emph{Graph Theory in Memory of G. A. Dirac}, Annals of Discrete Mathematics \textbf{41}, North-Holland, 1989, pp. 371--396.
\end{thebibliography}
\end{document}
